\scanpage{80}
% archive foliation: 79
% lines 1--3
Groupe profini \uline{spécial}

\deleted{\unsure{1}} C'est un système \inserted{$\Sigma=$}
$(T,\pi,I,(\Phi_i)_{i\in I}\in(\operatorname{Homext}(T,\pi))^I$
avec
% lines 4--7
\[\Sigma\left\{\begin{array}{l}
T\ \text{groupe profini}\ \simeq\widehat Z\\
\pi\ \text{groupe profini}\\
I\ \text{ensemble fini}\\
\Phi_i\in\operatorname{Homext}(T,\pi)\qquad\text{pour }i\in I
\end{array}\right.\]
% lines 8--10
avec l'exigence que ce système est
\deleted{\unsure{équivalent}}\inserted{isomorphe}
à l'un des systèmes $\Sigma_\nu$ ($\nu\in N$, $\nu\ge1$)
suivants :
% lines 11--15
\[\Sigma_\nu\left\{\begin{array}{l}
T=\widehat Z\\
\pi=\text{groupe profini}\ \inserted{\Pi_\nu=}\widehat L_\nu,\ \text{où }L_\nu=\left\{\lambda_0,\ldots,\lambda_\nu\mid\lambda_0\lambda_1\cdots\lambda_\nu\unsure{=1}\right.\\
\qquad(\text{donc }\pi\text{ prolibre à }\nu\text{ générateurs})\\
\unsure{\Phi|I}=[0,\nu]_N\\
\Phi_i=\text{classe de l'hom }\widehat Z\xrightarrow{\varphi_i}\pi\quad\text{tel que }\varphi_i(1)=\lambda_i
\end{array}\right.\]
% lines 16--22
\uline{NB} $I$ peut être interprété comme une \uline{partie} de $\operatorname{Homext}(T,\pi)$.
Si $\Sigma\simeq\Sigma_\nu$, $\nu$ est uniquement déterminé
comme $\operatorname{card}I-1$, ou \deleted{\unsure{peut}}
\deleted{\unsure{être}} comme rang profini de $\pi_{\mathrm{ab}}$.
Pour $\nu$ fixé, la catégorie des $\Sigma$ correspondants
est équivalente à la catégorie des torseurs
sous
\[\widetilde A_\nu=\operatorname{Aut}(\Sigma_\nu)\]
% lines 23--25
On a des homomorphismes
\[\begin{aligned}
\widetilde A_\nu&\longrightarrow\unsure{\mathfrak S_{\nu+1}|\mathcal S_{\nu+1}}\ (=\operatorname{Aut}(I_\nu))\\
\widetilde A_\nu&\xrightarrow{\chi}\operatorname{Aut}\widehat Z\simeq\widehat Z^*
\end{aligned}\]
